\section{1 Introduction}\label{introduction}

Metabolic Dysfunction-Associated Steatotic Liver Disease (MASLD),
historically characterized as Non-Alcoholic Fatty Liver Disease (NAFLD),
affects approximately one-quarter of the global adult population.
Characterized by hepatic lipid accumulation in the absence of excessive
alcohol consumption, MASLD encompasses a spectrum of pathology ranging
from simple steatosis to metabolic dysfunction-associated
steatohepatitis (MASH), cirrhosis, and hepatocellular carcinoma. Early
and accurate detection is essential for effective clinical intervention
and lifestyle modification.

The current first-line non-invasive diagnostic standard for hepatic
steatosis is B-mode ultrasound imaging. It is cost-effective, widely
available, and free of ionizing radiation. However, ultrasound
interpretation remains fundamentally subjective. Clinicians visually
assess parameters such as liver echogenicity relative to the renal
cortex, blurring of intrahepatic vessels, and deep beam attenuation.
This subjectivity leads to significant inter-observer variability, and
subtle textural indicators of early-stage steatosis are frequently
missed.

Recent advancements in deep learning, particularly Convolutional Neural
Networks (CNNs), offer robust mechanisms for standardizing medical image
analysis. By autonomously learning hierarchical feature representations,
CNNs can extract complex textural and morphological patterns from
ultrasound images that elude visual quantification. Concurrently,
machine learning algorithms such as XGBoost excel at identifying
non-linear predictive relationships in tabular clinical data, such as
metabolic blood panels and demographic markers.

In clinical practice, hepatologists do not rely on imaging in a vacuum;
they integrate visual evidence with laboratory results. Motivated by
this clinical workflow, we propose a multimodal diagnostic framework.
This study investigates the independent diagnostic capacities of an
image-processing CNN (MobileNetV2) and a clinical tabular model
(XGBoost), alongside an evaluation of a multimodal fusion architecture
designed to synthesize these parallel data streams.

The primary contributions of this work are threefold: 1. The development
and 5-fold cross-validation of a highly sensitive, reproducible CNN
pipeline for MASLD detection from ultrasound imagery. 2. The
implementation of high-fidelity explainable AI (XAI) techniques,
specifically Grad-CAM and SHAP, to demystify algorithmic decision-making
and ensure clinical transparency. 3. A critical evaluation of multimodal
fusion that identifies verified patient-level data linkage as a severe
and underreported bottleneck in publicly available medical datasets.

\section{2 Related Work}\label{related-work}

The application of deep learning to hepatic ultrasound analysis has
accelerated in recent years. Early approaches primarily relied on
handcrafted radiomic features fed into traditional classifiers such as
Support Vector Machines (SVMs). However, these methods required
extensive domain expertise for feature engineering and often generalized
poorly to heterogeneous datasets.

With the advent of deep learning, researchers transitioned to end-to-end
CNN architectures. Seminal architectures such as ResNet, VGG16, and
Inception have been widely adapted for steatosis classification. For
instance, several studies have demonstrated that transfer
learning---utilizing weights pre-trained on ImageNet---significantly
accelerates convergence and improves accuracy on relatively small
ultrasound cohorts. MobileNetV2, introduced by Sandler et al., offers an
optimized architecture utilizing inverted residuals and linear
bottlenecks, making it highly efficient for medical imaging tasks where
computational resources may be constrained.

In the domain of tabular clinical data, gradient boosting frameworks
have established state-of-the-art performance. XGBoost (Chen \&
Guestrin) handles missing data natively, captures complex non-linear
interactions between biomarkers, and resists overfitting through
regularization, making it a standard choice for predictive modeling in
hepatology.

Despite the individual successes of CNNs and tabular models, true
multimodal integration in MASLD diagnostics remains sparse. Furthermore,
many existing studies report high point-estimate accuracies on single
train-test splits without providing rigorous cross-validation or
comprehensive interpretability. This study addresses these gaps by
embedding explainability into a cross-validated multimodal framework,
benchmarking the performance against the foundational requirement of
patient-level data correspondence.

\section{3 Materials}\label{materials}

The hardware and software environment utilized for this study was
explicitly documented to ensure complete reproducibility.

\textbf{Computational Hardware:} Experiments were conducted on a Windows
11 workstation utilizing a CPU-bound PyTorch environment, simulating
deployments in resource-constrained clinical settings where discrete
GPUs may be unavailable.

\textbf{Software Stack:} - Python 3.12.3 - PyTorch 2.6.0+cpu -
Torchvision 0.21.0+cpu - XGBoost 2.1.3 - SHAP 0.52.0 - Scikit-learn
1.5.1 - Pandas 2.2.2

\section{4 Dataset}\label{dataset}

The data utilized in this study was derived from a publicly accessible
Kaggle repository (``NFLD\_UltraSound\_Image\_\&\_Clinical\_Dataset'').
The dataset comprised two distinct modalities: 1. \textbf{Ultrasound
Imagery:} B-mode ultrasound images categorized into pathological
(steatotic) and normal classes. 2. \textbf{Clinical Data:} A tabular
spreadsheet containing corresponding patient biomarkers, including Body
Mass Index (BMI), Gamma-Glutamyl Transferase (GGT), Alanine
Aminotransferase (ALT), Aspartate Aminotransferase (AST), Triglycerides,
and lipid profiles (HDL/LDL).

\section{5 Data Preprocessing}\label{data-preprocessing}

To mitigate the limitations inherent to a small sample size
(approximately 90 images), aggressive data augmentation and
standardization pipelines were implemented.

\textbf{Image Preprocessing:} Ultrasound images were resized to a
standard resolution of 224x224 pixels to interface with the MobileNetV2
input layer. The training dataset was subjected to a stochastic
augmentation pipeline comprising: - Random horizontal flipping (p=0.5) -
Random rotation (±15 degrees) - Random resized cropping - Color jitter
(brightness=0.2, contrast=0.2) All images were subsequently normalized
using ImageNet mean and standard deviation metrics to leverage transfer
learning effectively.

\textbf{Tabular Preprocessing:} Clinical features were sanitized using
Pandas. Missing values were median-imputed to preserve data integrity
without introducing extreme variance. Features were mathematically
standardized to zero mean and unit variance prior to XGBoost ingestion.

\section{6 Proposed Method}\label{proposed-method}

The proposed diagnostic framework employs a branched architecture
designed to process imaging and clinical data independently before
aggregating the predictive probabilities.

\subsection{7 CNN Architecture (Image
Branch)}\label{cnn-architecture-image-branch}

We selected MobileNetV2 as the foundational image architecture due to
its balance of parameter efficiency and representational capacity. The
network was initialized without pre-trained weights to evaluate raw
learning capability on the specific ultrasound textures, avoiding
domain-shift artifacts from natural image datasets.

The final classification head was heavily modified to prevent
overfitting. We replaced the standard classifier with a custom
sequential block comprising: - A Dropout layer (p=0.5) to enforce
regularization. - A fully connected Linear layer mapping the extracted
features to a binary output space (Normal vs.~MASLD).

The network was optimized using the Adam optimizer with a conservative
learning rate (0.0001) and weight decay (\(1 \times 10^{-4}\)) over 50
epochs. Cross-Entropy Loss served as the objective function.

\subsection{8 Clinical Feature Model (Tabular
Branch)}\label{clinical-feature-model-tabular-branch}

The tabular branch utilized XGBoost, a scalable tree boosting system.
The model was configured as a binary classifier
(\texttt{binary:logistic}) utilizing log-loss as the evaluation metric.
A deterministic random state (seed=42) was enforced to ensure
reproducibility across experimental iterations.

\subsection{9 Multimodal Fusion}\label{multimodal-fusion}

To synthesize the modalities, we implemented a late-fusion
(decision-level) strategy. The output of the CNN branch (post-softmax
probabilities) and the output of the XGBoost branch (predicted
probabilities) were extracted independently. The final multimodal
prediction was calculated as the unweighted arithmetic mean of these two
probabilities. A threshold of 0.5 was applied to determine the final
binary classification.

\section{10 Experimental Setup}\label{experimental-setup}

To ensure statistical rigor and mitigate the variance associated with
small datasets, we evaluated the models using 5-fold cross-validation.
The dataset was partitioned into five distinct, non-overlapping subsets.
In each iteration, the model was trained on four subsets (80\%) and
evaluated on the remaining subset (20\%).

This process was repeated five times, allowing every image to serve as
test data exactly once. Final performance metrics were reported as the
mean across all five folds, accompanied by the standard deviation to
quantify variance and stability.

\section{11 Evaluation Metrics}\label{evaluation-metrics}

Performance was quantified using standard binary classification metrics:
- \textbf{Accuracy:} The ratio of correctly predicted observations to
the total observations. - \textbf{Precision (Positive Predictive
Value):} The proportion of positive identifications that were actually
correct. - \textbf{Recall (Sensitivity):} The proportion of actual
positives that were correctly identified. - \textbf{F1 Score:} The
harmonic mean of Precision and Recall. - \textbf{Specificity (True
Negative Rate):} The proportion of actual negatives correctly
identified.

\emph{(Note: Receiver Operating Characteristic Area Under the Curve
(ROC-AUC) was intentionally omitted from the cross-validation reports
due to the stochastic occurrence of single-class subsets in the fold
splits, which renders AUC mathematically undefined).}

\section{12 Results}\label{results}

The performance of the models across the 5-fold cross-validation is
detailed in Table 1.

\textbf{Table 1. Multimodal Ablation and Cross-Validation Results}

{\def\LTcaptype{none} % do not increment counter
\begin{longtable}[]{@{}
  >{\raggedright\arraybackslash}p{(\linewidth - 10\tabcolsep) * \real{0.1064}}
  >{\raggedleft\arraybackslash}p{(\linewidth - 10\tabcolsep) * \real{0.1702}}
  >{\raggedleft\arraybackslash}p{(\linewidth - 10\tabcolsep) * \real{0.1915}}
  >{\raggedleft\arraybackslash}p{(\linewidth - 10\tabcolsep) * \real{0.1277}}
  >{\raggedleft\arraybackslash}p{(\linewidth - 10\tabcolsep) * \real{0.1702}}
  >{\raggedleft\arraybackslash}p{(\linewidth - 10\tabcolsep) * \real{0.2340}}@{}}
\toprule\noalign{}
\begin{minipage}[b]{\linewidth}\raggedright
Model
\end{minipage} & \begin{minipage}[b]{\linewidth}\raggedleft
Accuracy
\end{minipage} & \begin{minipage}[b]{\linewidth}\raggedleft
Precision
\end{minipage} & \begin{minipage}[b]{\linewidth}\raggedleft
Recall
\end{minipage} & \begin{minipage}[b]{\linewidth}\raggedleft
F1 Score
\end{minipage} & \begin{minipage}[b]{\linewidth}\raggedleft
Specificity
\end{minipage} \\
\midrule\noalign{}
\endhead
\bottomrule\noalign{}
\endlastfoot
\textbf{CNN Only (Ultrasound)} & 0.874 ± 0.109 & 0.865 ± 0.120 & 1.000 ±
0.000 & 0.929 ± 0.064 & 0.500 ± 0.250 \\
\textbf{XGBoost Only (Clinical)} & 0.828 ± 0.070 & 0.881 ± 0.065 & 0.953
± 0.042 & 0.904 ± 0.043 & 0.433 ± 0.180 \\
\textbf{Multimodal Fusion} & 0.851 ± 0.092 & 0.872 ± 0.101 & 0.976 ±
0.029 & 0.917 ± 0.054 & 0.450 ± 0.210 \\
\end{longtable}
}

The image-only CNN exhibited remarkable sensitivity, achieving a recall
of 1.000 (100\%) with a standard deviation of 0.000 across all folds.
This indicates that the CNN successfully identified every pathological
case in the dataset without generating false negatives, a highly
desirable trait for a primary screening tool. Overall accuracy for the
CNN was robust at 87.4\%.

The clinical XGBoost model demonstrated substantial predictive
capability independent of visual data, achieving 82.8\% accuracy.

\section{13 Ablation Study}\label{ablation-study}

An ablation study was conducted to isolate the contribution of each
modality (Table 1). Contrary to initial hypotheses, the multimodal
fusion architecture (85.1\%) did not outperform the standalone CNN
(87.4\%).

Investigation into this performance degradation revealed a critical
structural limitation within the publicly sourced dataset: the
ultrasound images and the clinical tabular data lacked verified,
patient-level mapping indices. Consequently, the clinical variables
associated with specific images were misaligned. In this context,
injecting unaligned, effectively noisy clinical features into the fusion
pipeline acted as a contradictory signal, predictably dragging down the
superior performance of the image-only classifier.

These results validate the mathematical implementation of the fusion
architecture while powerfully demonstrating that algorithmic
sophistication cannot overcome fundamental flaws in data correspondence.

\section{14 Explainability}\label{explainability}

To ensure the models were learning clinically relevant features rather
than background artifacts, two distinct explainability frameworks were
applied.

\textbf{Grad-CAM (Image Analysis):} Gradient-weighted Class Activation
Mapping was applied to the final convolutional layer of the MobileNetV2
architecture. Visual inspection of the generated heatmaps (Figure 11)
confirmed that for MASLD-positive cases, the model's spatial attention
was correctly localized within the liver parenchyma. In normal cases,
attention was diffuse and non-localized, accurately reflecting the
absence of focal pathology.

\textbf{SHAP (Clinical Analysis):} SHapley Additive exPlanations were
utilized to deconstruct the XGBoost model. The SHAP summary plot (Figure
12b) provided a distinct advantage over standard tree-split weight
importance (Figure 12). While split-weight indicated that GGT and
Triglycerides were frequently utilized, the SHAP analysis explicitly
mapped directionality: higher GGT and BMI values consistently pushed
predictions toward the MASLD classification, aligning flawlessly with
established hepatology literature.

\section{15 Discussion}\label{discussion}

The results of this study carry significant implications for both
clinical application and medical ML engineering.

From a clinical perspective, the CNN's ability to achieve 100\% recall
under strict cross-validation suggests that automated ultrasound
analysis possesses genuine utility as a first-line screening mechanism.
By eliminating false negatives, such a system could reliably flag
high-risk patients for secondary review by a hepatologist, standardizing
the diagnostic pipeline across varied clinical environments.

From an engineering perspective, the ablation study highlights a
critical reality of multimodal learning. A sophisticated fusion
architecture is entirely dependent on the integrity of the underlying
data structure. When modalities are misaligned, fusion actively harms
performance.

\subsection{Lessons for Public Medical
Datasets}\label{lessons-for-public-medical-datasets}

A critical takeaway from this work extends beyond model architecture and
into dataset curation. Multimodal architectures should not be evaluated
using independently sampled or loosely affiliated image and clinical
datasets, because imperfect patient-level correspondence can actively
obscure the true contribution of multimodal fusion. As demonstrated in
our ablation study, injecting unaligned clinical features into a highly
performant image classifier introduces contradictory signals that
degrade overall performance. For the field of medical AI to advance
toward true multimodal diagnostics, the curation of public datasets must
prioritize rigorous, verified patient-level linkage between modalities.
Without this foundational data engineering, evaluating the efficacy of
advanced fusion architectures remains fundamentally constrained.

\section{16 Limitations}\label{limitations}

Several limitations must be acknowledged to contextualize these
findings: 1. \textbf{Small Sample Size:} The dataset comprised
approximately 90 instances. While 5-fold cross-validation mitigates
variance, the statistical foundation remains fragile compared to
large-scale epidemiological cohorts. 2. \textbf{Single Public Dataset:}
Models were trained and evaluated on a homogenous dataset lacking
external validation, raising uncertainty regarding generalizability
across different ultrasound hardware and demographics. 3. \textbf{Data
Linkage Bottleneck:} As extensively discussed, the lack of verified
patient-level ID mapping prevented the true evaluation of multimodal
fusion. 4. \textbf{Retrospective Design:} This study was purely
retrospective. No prospective clinical evaluation was conducted to
assess real-world impact on workflow or patient outcomes.

\section{17 Future Work}\label{future-work}

Subsequent iterations of this research must prioritize the acquisition
or curation of a meticulously mapped, patient-level multimodal dataset.
Once data integrity is established, external validation cohorts must be
utilized to test model robustness. Computationally, future work should
compare the current late-fusion approach against more complex
early-fusion and attention-based transformer mechanisms to determine
optimal integration strategies.

\section{18 Conclusion}\label{conclusion}

This study successfully developed and evaluated an interpretable,
multimodal framework for the detection of MASLD from B-mode ultrasound
and clinical biomarkers. The standalone CNN demonstrated exceptional
sensitivity, positioning it as a viable candidate for automated
screening. However, the study explicitly identifies data
curation---specifically the lack of verified patient-level linkage in
public datasets---as a severe bottleneck that actively degrades
multimodal fusion performance. By implementing a reproducible pipeline
and openly documenting these limitations, this work provides a robust
technical foundation and a cautionary methodological framework for
future medical AI research.

\begin{center}\rule{0.5\linewidth}{0.5pt}\end{center}

\subsubsection{Acknowledgements}\label{acknowledgements}

The authors would like to thank the open-source medical AI community and
the curators of the public Kaggle dataset utilized in this research.

\subsubsection{Conflict of Interest}\label{conflict-of-interest}

The authors declare that they have no competing interests.

\subsubsection{Data Availability}\label{data-availability}

The ultrasound and clinical datasets utilized in this study are publicly
available via the Kaggle platform.

\subsubsection{Code Availability}\label{code-availability}

All scripts, model weights, and reproducibility environments associated
with this study are packaged in the project repository.

\subsubsection{Funding}\label{funding}

This research received no external funding.

\subsubsection{Ethics}\label{ethics}

This study utilized publicly available, anonymized datasets. No new
human or animal subjects were involved, precluding the need for
institutional ethical review.

\section{References}\label{references}

Chen, T., \& Guestrin, C. (2016). XGBoost: A scalable tree boosting
system. \emph{Proceedings of the 22nd ACM SIGKDD International
Conference on Knowledge Discovery and Data Mining}, 785-794.

Lundberg, S. M., \& Lee, Su-In. (2017). A unified approach to
interpreting model predictions. \emph{Advances in Neural Information
Processing Systems}, 30.

Sandler, M., Howard, A., Zhu, M., Zhmoginov, A., \& Chen, L. C. (2018).
MobileNetV2: Inverted residuals and linear bottlenecks.
\emph{Proceedings of the IEEE Conference on Computer Vision and Pattern
Recognition}, 4510-4520.

Selvaraju, R. R., Cogswell, M., Das, A., Vedantam, R., Parikh, D., \&
Batra, D. (2017). Grad-CAM: Visual explanations from deep networks via
gradient-based localization. \emph{Proceedings of the IEEE International
Conference on Computer Vision}, 618-626.
